\selectlanguage{english}
\noindent\textbf{Chapter I, first part: from presheaf operations to homotopical base change}

These pages construct a language in which a homotopy theory can carry its own diagram categories. The point of departure is elementary: a functor of indexing categories gives restriction of presheaves, restriction has Kan-extension adjoints when the required limits exist, and special inclusions admit further adjoints when the coefficient category is pointed. What changes in passing to homotopy theory is the meaning of a diagram. A diagram in a homotopy category records only the images of objects and arrows after localization. Localizing a category of diagrams retains additional morphisms and coherence. Grothendieck proposes to organize all these localized categories together, so that a universal construction can be performed on a coherent diagram and remain compatible with further parameters.

The historical motivation is stated explicitly in \emph{Pursuing Stacks}, \S69: a single derived category does not provide an adequate setting for the trace and determinant identities encountered around SGA~5 \citep[\S69]{Grothendieck1983}. Categories of filtered complexes and diagrams of complexes supply more information, but they too must be organized into a structure stable under introducing new diagram shapes. The letter to Thomason of 2 April 1991 describes the renewed investigation of this intrinsic structure \citep[pp.~2--4]{Grothendieck1991Letter}. Heller's homotopy theories form a closely related independent axiomatic development \citep{Heller1988}. The present opening chapter is exploratory: proposed axioms are changed, conclusions are withdrawn, and the final definitions are reached through computations that reveal which hypotheses are actually needed.

\medskip\noindent\textbf{The convention that governs the calculations.}
We work in a fixed universe of small indexing categories; a restricted domain must contain every opposite, product, coproduct, and comma category used in a particular assertion. The manuscript uses presheaves,
\[
 \DD_A(I)=\operatorname{Fun}(I^{\mathrm{op}},A),
 \qquad f_!\dashv f^*\dashv f_*.
\]
Thus an object over $[1]=(0\to1)$ has an underlying arrow $F_1\to F_0$. Restriction reverses both functors and natural transformations, so the explicit modern source is $\Dia^{\mathrm{coop}}$. The manuscript calls a theory with right Kan extensions a derivator \emph{\`a gauche}, because these are the operations of left exactness. A theory with left Kan extensions is called \emph{\`a droite}. The systematic axioms starting at manuscript p.~41 replace the earlier provisional numbering. These changes of convention are substantive reading instructions: a remembered modern formula with the opposite variance can reverse the cofinality condition.

The ordinary pointwise formulas already display the relevant slices:
\[
 (f_*F)_y=\lim_{(f/y)^{\mathrm{op}}}F,
 \qquad
 (f_!F)_y=\operatorname*{colim}_{(y/f)^{\mathrm{op}}}F.
\]
For derivators these are homotopy Kan formulas, expressed by invertibility of specified base-change maps. The map, not merely an isomorphism between its source and target, is part of the assertion. This explains the care devoted to adjunction units, counits, and the two equivalent constructions of the base-change mate.

\medskip\noindent\textbf{Supports, coherent squares, and shifts.}
Open inclusions are sieves and closed inclusions are cosieves. In the presheaf convention a closed inclusion extends by terminal objects under $i_*$; an open inclusion extends by initial objects under $j_!$. Pointedness identifies the two values with zero and permits the exceptional adjoints $i^!$ and $j^?$. The commentary proves these assertions from their universal properties and works out the nonsplit short exact sequence of arrow diagrams underlying the support decomposition. It also gives the finite-ordinal adjunction chain, including its four exceptional endpoints, while retaining the edition's correction of the stronger claim about every monotone map.

The same nonsplit extension exhibits a nonzero morphism between coherent derived diagrams whose components are all zero. Consequently the underlying-diagram functor is conservative without being faithful. This is precisely the distinction the localization axiom requires. Homotopy cartesian squares are coherent squares obtained by right Kan extension from a missing-corner diagram; they need not be ordinary pullbacks in the base homotopy category. Strongness, when imposed, addresses lifting arrows and morphisms of arrows. It is a separate condition from conservativity.

The shift $\DD_I(J)=\DD(I\times J)$ turns an entire coherent $I$-diagram into one coefficient. Its elementary formal properties explain the repeated demand that each new axiom survive shifting. The completed shift theorem and product base change establish the proposed morphisms induced by $f^*$, $f_!$, and $f_*$. At manuscript p.~18 the question whether evaluation at the base category is fully faithful also has a concrete negative answer for unrestricted morphisms: the invisible extension class yields a nonzero modification between morphisms out of the represented derivator of $[1]$ whose component at the point is zero. The example does not impose continuity on these morphisms.

For a Quillen model category, existence of the homotopy Kan extensions is a theorem rather than a formal property of localization. Cisinski supplies a precise general treatment \citep[Theorem~6.11]{Cisinski2003}. In the $\infty$-categorical formulation one takes
\[
 \DD_{\mathcal C}(I)=h\operatorname{Fun}(N(I)^{\mathrm{op}},\mathcal C);
\]
the pointwise extensions of a complete and cocomplete $\infty$-category give the derivator operations \citep[\S5.3]{RiehlVerity2017}. Reconstruction results for the whole prederivator retain their hypotheses on size and morphisms. Renaudin's comparison concerns combinatorial model theories and adjunctions \citep[Theorems~3.3.2 and~3.4.4]{Renaudin2009}, while Arlin distinguishes strict transformations from pseudonatural ones \citep[Theorems~19 and~29]{Arlin2020}. These are stronger invariants than the base category alone.

\medskip\noindent\textbf{The exactness question has a finiteness boundary.}
The connective derived theory $\mathsf{HOTAB}$ is additive and pointed but is not stable. Its abelianization morphism from spaces preserves homotopy colimits. It fails to preserve homotopy pullbacks: for two discrete two-point spaces over a point, the comparison is a map from $\mathbb Z^4$ to $\mathbb Z^3$ with an explicitly displayed rank-one kernel. This separates the true pushout assertion from the false pullback assertion in the same passage.

Stability answers the later ``Perplexit\'e'' after the scope of exactness is fixed. Between stable derivators, preservation of zero objects and cartesian squares is equivalent to preservation of zero objects and cocartesian squares \citep[Corollary~4.17]{Groth2013}. Homotopy finite limits commute with homotopy colimits in the corresponding precise stable theorem \citep[Theorem~2.14]{GrothShulman2017}. Arbitrary products and coproducts still need not commute. The diagonal array of copies of a field gives the comparison $k^{(\mathbb N)}\to k^{\mathbb N}$, which omits every sequence with infinite support. The modern finite theorem therefore does not validate an unrestricted preservation claim for all Kan extensions.

\medskip\noindent\textbf{Cofinality and universal base change.}
The decisive correction in the chapter concerns the coefficient system. Contractibility of every $f/y$ does not make restriction fully faithful on all coefficients. For the inclusion of the initial endpoint into $[1]$, every such comma category is a point, but the unit on $F_1\to F_0$ contains the original structure map. The example $0\to\mathbb Z$ disproves the proposed axiom. The surviving proof of the later cancelled corollary contains a valid slice identification; its last inference uses exactly the false, opposite-orientation cofinality clause. The audit identifies this inference without inventing an additional gap in the calculation.

For right direct image in the manuscript's convention, the correct cofinality test uses arrows from an object of the target to the image of the functor. Applied pointwise to a square, it produces categories of factorizations of a fixed base arrow. Their weak contractibility is the universal homotopy-exactness criterion. Properness specializes it to the comparison between a strict fibre and a comma category; smoothness is the dual incoming-lift condition. The commentary completes the composition proof by identifying the relevant lift-category functor as a strict pullback and applying Quillen's Theorem~A \citep[\S1]{Quillen1973}.

The conjectured converse is established: spaces alone detect universal homotopy exactness, so the proposed simultaneous use of $\mathsf{HOT}$ and $\mathsf{HOTAB}$ is unnecessary \citep[Theorem~3.16]{GrothPontoShulman2014}. Proper and smooth functors are treated systematically by Maltsiniotis \citep[Theorem~1.2]{Maltsiniotis2005Asphericite}. Strict pullbacks matter here. An equivalence of categories that misses an object can fail both properness and smoothness, as the two-object groupoid example shows.

The commentators' marked cone construction makes the strength of the homotopical test explicit. Adjoin an initial object to a small category $K$, send that object to $0$, and send $K$ to $1$. The resulting functor to $[1]$ is always smooth and a weak equivalence of nerves; it is proper exactly when $NK$ is weakly contractible. If $K$ is nonempty and connected, all its strict base-change squares pass the right Kan test for every complete ordinary coefficient category. A four-object crown gives $|NK|\simeq S^1$ and a failure for space coefficients, detected by the degree components of $\operatorname{Map}(S^1,S^1)$. The construction, including its universal ordinary-coefficient assertion, is proved in full and carries no claim of priority.

\medskip\noindent\textbf{The final definition and its scope.}
The corrected notion of $\DD$-equivalence compares every morphism between constant coherent coefficients. By Yoneda it is equivalent, when the representing objects exist, to invariance of constant-coefficient homology or cohomology. Its two-out-of-three property is therefore already formal at the prederivator level. The local versions distinguish discrete fibrations from discrete opfibrations and lead to asphericity and coasphericity. The omitted stability check under all local isomorphisms is supplied by an explicit isomorphism of comma categories. Locally constant coefficients are handled by a canonical morphism on a slice and conservativity.

The last proposition turns this constant-coefficient test into a criterion for arbitrary coefficients by varying the target in a Yoneda argument. Its Hom-set formulation needs left Kan extensions only; its cohomological formulation additionally needs the absolute right Kan extensions that define $H^*$. This is an existence qualification, not a failure of the argument. The questions treated here consequently have precise answers once variance, coefficient class, finiteness, and existence of adjoints are stated. The attribution to Thomason, which the text leaves unspecified, remains a documentary uncertainty.
